\documentclass[aspectratio=169,11pt]{beamer}

\usetheme{Madrid}
\usecolortheme{default}
\usefonttheme{professionalfonts}
\setbeamertemplate{navigation symbols}{}
\setbeamertemplate{footline}[frame number]

\usepackage{amsmath, amssymb, booktabs, graphicx}

\title{Why Institutions Matter for Labor Markets}
\subtitle{Labor Markets and Political/Cultural Institutions --- Week 1}
\author{Teaching deck draft}
\date{}

\begin{document}

\begin{frame}
  \titlepage
\end{frame}

\begin{frame}{Central question}
\begin{itemize}
    \item What counts as a labor-market institution?
    \item Through which mechanisms do institutions shape labor allocation and inequality?
    \item How can labor economists study formal rules and informal norms without losing the labor object?
\end{itemize}
\end{frame}

\begin{frame}{Course framing}
\begin{itemize}
    \item The course starts from labor outcomes: wages, employment, informality, mobility, bargaining, sorting, and worker welfare.
    \item Institutions are not background context; they shape feasible contracts, enforcement, search, referrals, compliance, and political demand.
    \item The discipline: labor object, institutional channel, mechanism, evidence, and counterfactual.
\end{itemize}
\end{frame}

\begin{frame}{Formal and informal institutions}
\small
\begin{tabular}{p{0.26\linewidth} p{0.31\linewidth} p{0.33\linewidth}}
\toprule
Type & Examples & Labor-market relevance \\
\midrule
Formal & law, contracts, regulation, unions, enforcement, state capacity & legal constraints, bargaining, incidence, compliance \\
Informal & culture, norms, trust, family ties, networks, hierarchy, status beliefs & search, mobility, referrals, sorting, compliance, regulation demand \\
Hybrid & social enforcement of rules, legitimacy, informal compliance & whether formal rules bind in practice \\
\bottomrule
\end{tabular}
\end{frame}

\begin{frame}{Objects, mechanisms, outcomes}
\[
Y_{ifrt}=g(O_{rt},M_{ifrt},X_{ifrt})
\]
\begin{itemize}
    \item Object: law, contract regime, norm, network, hierarchy, family obligation, or belief system.
    \item Mechanism: incentives, bargaining, enforcement, information, trust, insurance, sorting, mobility, compliance.
    \item Outcome: wages, employment, hours, informality, migration, job access, contract form, inequality.
\end{itemize}
\end{frame}

\begin{frame}{Mechanism map}
\small
\begin{tabular}{p{0.24\linewidth} p{0.31\linewidth} p{0.31\linewidth}}
\toprule
Mechanism & Formal side & Informal side \\
\midrule
Incentives & taxes, benefits, legal constraints & work norms, reciprocity, aspirations \\
Bargaining & dismissal rules, bargaining law & authority norms, status hierarchy \\
Compliance & inspections, courts, penalties & legitimacy, shame, network sanctions \\
Information & disclosure, contract rules & referrals, social learning, inherited beliefs \\
Insurance & UI, severance, social insurance & kinship, caste, community support \\
Mobility and sorting & licensing, contract boundaries & hierarchy, network lock-in, local status \\
\bottomrule
\end{tabular}
\end{frame}

\begin{frame}{Comparative measurement}
\begin{itemize}
    \item Freeman: labor-market institutions differ systematically across economies.
    \item Botero et al.: formal labor regulation can be coded and compared.
    \item Measurement is not the mechanism: an index must still be linked to wages, employment, adjustment, or inequality.
    \item Cross-country designs need caution about bundled institutions and omitted development paths.
\end{itemize}
\end{frame}

\begin{frame}{Labor law, contracts, enforcement}
\begin{itemize}
    \item MacLeod: law and employment contracts shape what promises can be made credible.
    \item Statutes become labor-market constraints through courts, inspections, legal access, and administrative reach.
    \item The effective rule is not only what is written; it is what workers and firms expect to be enforced.
    \item This is the bridge to Week 2.
\end{itemize}
\end{frame}

\begin{frame}{Culture as informal institution}
\begin{itemize}
    \item Culture includes beliefs about work, authority, mobility, unemployment, family obligation, status, trust, and regulation.
    \item Gender norms are one application, not the whole domain.
    \item Informal institutions affect labor supply, migration, referrals, compliance, bargaining, occupational allocation, and support for regulation.
    \item Guiso, Sapienza, and Zingales provide the broad culture-as-mechanism frame.
\end{itemize}
\end{frame}

\begin{frame}{Family values and labor regulation}
\begin{itemize}
    \item Alesina, Algan, Cahuc, and Giuliano link family values to demand for labor regulation.
    \item Informal insurance and family support can substitute for or complement formal worker protection.
    \item Political support for regulation is itself a labor-market channel.
    \item Culture can shape formal rules, not only individual choices.
\end{itemize}
\end{frame}

\begin{frame}{Networks, mobility, misallocation}
\begin{itemize}
    \item Munshi and Rosenzweig: caste and community networks can insure workers and transmit job information.
    \item The same networks can discourage migration or occupational switching when exit weakens insurance.
    \item A rural-urban wage gap is not enough to predict mobility when informal institutions bind.
    \item Network benefits and network lock-in must be studied together.
\end{itemize}
\end{frame}

\begin{frame}{Why informal institutions persist}
\begin{itemize}
    \item Families transmit beliefs about work, risk, status, care, and mobility.
    \item Networks create referrals and insurance that make exit costly.
    \item Employers reuse trusted channels because they reduce hiring uncertainty.
    \item Social sanctions and political preferences can outlast the formal rule that first created them.
\end{itemize}
\end{frame}

\begin{frame}{Empirical strategies}
\small
\begin{tabular}{p{0.30\linewidth} p{0.32\linewidth} p{0.25\linewidth}}
\toprule
Strategy & Institutional object & Labor outcome \\
\midrule
Comparative measures & regulation, bargaining, enforcement & employment, inequality, adjustment \\
Legal variation & courts, inspections, contract rules & take-up, formality, wages \\
Historical persistence & norms, hierarchy, past institutions & mobility, sorting, regulation demand \\
Ancestry designs & inherited beliefs & labor supply, occupation, fertility \\
Network designs & referrals, insurance, community ties & migration, matching, wage gaps \\
Admin data & contracts, enforcement, workplace rules & incidence, compliance, sorting \\
\bottomrule
\end{tabular}
\end{frame}

\begin{frame}{Week 1 lab}
\begin{itemize}
    \item Reproduce: a synthetic inherited-norms and labor-outcomes factbook.
    \item Diagnose: classify each variable as institutional object, mechanism, or outcome.
    \item Transfer: apply the same logic to network insurance, migration, and wage gaps.
    \item The point is research design discipline, not official replication.
\end{itemize}
\end{frame}

\begin{frame}{Bridge to Week 2}
\begin{itemize}
    \item Week 1 gives the full institutional map.
    \item Week 2 zooms in on labor law, employment contracts, enforcement, and state capacity.
    \item Central next question: when do statutes become effective workplace constraints?
    \item Later weeks return to culture, norms, hierarchy, persistence, and reform.
\end{itemize}
\end{frame}

\begin{frame}{Takeaways}
\begin{itemize}
    \item Labor-market institutions include both formal rules and informal social structures.
    \item The same institution can change incentives, bargaining, compliance, information, insurance, sorting, and mobility.
    \item Culture is measurable and labor-relevant, but it must be connected to concrete mechanisms.
    \item Strong institutional labor research separates objects, mechanisms, outcomes, and counterfactuals.
\end{itemize}
\end{frame}

\end{document}
